\makeatletter
\def\input@path{{../styles/}{../../styles/}{../../../styles/}{../}{../../}{../../../}}
\makeatother
\documentclass{ee122_notes}
\input{macros.tex}
\input{preamble.tex}
\renewcommand{\releasedate}{April 1, 2026}


\begin{document}
\section*{EE 122/ME 141 Week 10, Lecture 2: Practice state-feedback control design}
\subsection*{Instructor: \instructor}
\subsection*{Date: \releasedate}
\section{Announcements}
\begin{itemize}
    \item Project problem statement out now. First milestone due April 12, 2026.
\end{itemize}
\section{Learning Objective}
\begin{enumerate}
    \item Practice the design of state-feedback controllers for 1-D and 2-D systems.
    \item Connect physical closed-loop behavior to desired poles.
    \item Understand the role of the feedforward gain \(k_f\) when the goal is reference tracking instead of only regulation.
\end{enumerate}

In this lecture, we will walk through three examples of state-feedback control design. The main theme in the lecture is that the physics of the system guides our choice of desired poles, which then leads us to the control laws. 

\section{Example 0: first-order step response}

Before starting the examples, we review why the phrase ``reach \(63\%\) of the final value in one time constant'' appears so often in first-order systems. Consider a stable first-order closed-loop transfer function written as
\[
G_{cl}(s)=\frac{b}{s+a},
\qquad a>0.
\]
If the reference input is a unit step,
\[
R(s)=\frac{1}{s},
\]
then the output is
\[
Y(s)=G_{cl}(s)R(s)=\frac{b}{s(s+a)}.
\]
Using partial fractions, we can write that \[
\frac{b}{s(s+a)}=\frac{A}{s}+\frac{B}{s+a}
\]
which implies
\[
b=A(s+a)+Bs.
\]
We get that $A=\frac{b}{a}$ and $B=-\frac{b}{a}$. So,
\[
Y(s)=\frac{b}{a}\left(\frac{1}{s}-\frac{1}{s+a}\right).
\]
Taking the inverse Laplace transform,
\[
y(t)=\frac{b}{a}\left(1-e^{-at}\right),
\]
and the final value is
\[
y(\infty)=\frac{b}{a},
\]
so the response can be written as
\[
y(t)=y(\infty)\left(1-e^{-at}\right).
\]
At the time
\[
t=\frac{1}{a},
\]
we get
\[
y\left(\frac{1}{a}\right)=y(\infty)\left(1-e^{-1}\right)
\approx 0.63\,y(\infty).
\]
This means that for a first-order system with pole at \(-a\), the time constant is \(1/a\), and the output reaches about \(63\%\) of its final value after one time constant.

\begin{popquiz}
If we want a first-order closed-loop response to reach \(63\%\) of its final value in \(1\) second, what should the closed-loop pole be?
\popqsplit

Since the time constant must be \(1\), we need
\[
\frac{1}{a}=1
\quad\Longrightarrow\quad
a=1.
\]
Therefore, the desired pole is
\[
s=-1.
\]
\end{popquiz}

\section{Example 1: first-order system control design}

Consider the first-order system
\[
\dot{x}=-0.5x+u,
\qquad
y=2x.
\]
We use this example to practice both state-feedback design and the derivation of the feedforward gain \(k_f\) for reference tracking.

\subsection*{Step 1: write the state-space matrices}

From the model, we can directly read off
\[
A=[-0.5],
\qquad
B=[1],
\qquad
C=[2],
\qquad
D=[0].
\]

\subsection*{Step 2: identify the open-loop pole}

For this scalar system, the open-loop pole is simply the eigenvalue of \(A\):
\[
\lambda_{ol}=-0.5.
\]
So the open-loop system is stable, but it is slower than what we want.

\subsection*{Step 3: the system requirement}

We want the closed-loop output \(y(t)\) to reach \(63\%\) of its final value in \(1\) second when a step reference is applied. From the popquiz above, we see that the desired closed-loop pole is
\[
\lambda_{cl}=-1.
\]

\subsection*{Step 4: choose the control law}

Since we want to track a reference, we use
\[
u=-Kx+k_f r.
\]
Substituting into the plant gives
\[
\dot{x}=-0.5x+\left(-Kx+k_f r\right).
\]
Therefore,
\[
\dot{x}=(-0.5-K)x+k_f r.
\]
The closed-loop pole is
\(
-0.5-K.
\)
To place this pole at \(-1\), we solve
\(
-0.5-K=-1.
\)
Hence,
\[
K=0.5.
\]

\subsection*{Step 5: derive the feedforward gain $k_f$}

Now we derive \(k_f\) step by step rather than quoting a formula.At steady state for a constant reference \(r\), we have
\(
\dot{x}=0.
\)
So the closed-loop state equation becomes
\[
0=(-0.5-K)x_{ss}+k_f r.
\]
The output is
\[
y=2x,
\]
so perfect tracking requires
\[
y_{ss}=r.
\]
Since \(y_{ss}=2x_{ss}\), this means
\[
2x_{ss}=r
\quad\Longrightarrow\quad
x_{ss}=\frac{r}{2}.
\]
Substitute this into the steady-state equation:
\[
0=(-0.5-K)\frac{r}{2}+k_f r.
\]
Assuming \(r\neq 0\), divide by \(r\):
\[
0=-\frac{0.5+K}{2}+k_f.
\]
Therefore,
\[
k_f=\frac{0.5+K}{2}.
\]
Using \(K=0.5\), we get
\[
k_f=\frac{0.5+0.5}{2}=\frac{1}{2}=0.5.
\]
So the controller is
\[
u=-0.5x+0.5r.
\]

\subsection*{Step 6: write the closed-loop system and transfer function}

Substituting \(K=0.5\) and \(k_f=0.5\) into the closed-loop equation gives
\[
\dot{x}=-x+0.5r,
\qquad
y=2x.
\]
Taking Laplace transforms with zero initial conditions,
\[
sX(s)=-X(s)+0.5R(s).
\]
So
\[
(s+1)X(s)=0.5R(s),
\qquad
X(s)=\frac{0.5}{s+1}R(s).
\]
Since \(y=2x\),
\[
Y(s)=2X(s)=\frac{1}{s+1}R(s).
\]
Therefore, the closed-loop transfer function from \(r\) to \(y\) is
\[
\frac{Y(s)}{R(s)}=\frac{1}{s+1}.
\]
This is the first-order form we wanted, with pole at \(-1\)!! All done!

\subsection*{Step 7: sanity check}

For a unit step reference,
\[
R(s)=\frac{1}{s}.
\]
Hence,
\[
Y(s)=\frac{1}{s(s+1)}.
\]
Using partial fractions,
\[
\frac{1}{s(s+1)}=\frac{1}{s}-\frac{1}{s+1},
\]
and taking inverse Laplace transform, we get
\[
y(t)=1-e^{-t}.
\]
At \(t=1\),
\[
y(1)=1-e^{-1}\approx 0.63.
\]
So the output reaches about \(63\%\) of its final value at \(1\) second, as required.

\section{Example 2: deriving the reachable/controllable canonical form}

Consider the DC motor transfer function
\[
G(s)=\frac{2}{2s+1}.
\]
We first rewrite it in a form that makes the pole easier to see:
\[
G(s)=\frac{1}{s+0.5}.
\]

\subsection*{Step 1: write a state-space model in reachable canonical form}
Use the approach we discussed in class: multiply numerator and denominator by an unknown function \(Z(s)\):
\[
G(s)=\frac{1}{s+0.5}=\frac{Z(s
)}{Z(s)(s+0.5)}.
\]
Then, we write $Y(s) = 1 \cdot Z(s)$ and $U(s) = Z(s)(s+0.5)$, which gives
\[
Y(s)=Z(s),
\qquad
U(s)=Z(s)(s+0.5).
\]
Now, let's define the state as
\[
x=z
\]
so we get,
\[
y = x,
\qquad
u = \dot x + 0.5 x.
\]
Rearranging gives
\[
\dot x = -0.5 x + u.
\]
Thus, a state-space representation in reachable canonical form is
\[
A=[-0.5],
\qquad
B=[1],
\qquad
C=[1],
\qquad
D=[0].
\]

\subsection*{Step 2: identify the open-loop pole}

Again, for this scalar system, the open-loop pole is
\[
\lambda_{ol}=-0.5.
\]

\subsection*{Step 3: convert the physical requirement into a desired pole}

We want the closed-loop output to reach \(63\%\) of its final value in \(1\) second under a step input. Using the first-order derivation from the beginning of the lecture, the desired pole is
\[
\lambda_{cl}=-1.
\]


\subsection*{Step 4: compute the state-feedback gain}
We get 
$K=0.5$, same as before. 

\section{Practice problem (take home): second-order system control design}

Consider the second-order system
\[
\dot{x}=Ax+Bu,
\qquad
y=Cx+Du
\]
with
\[
A=\begin{bmatrix}
-2 & -13\\
1 & 0
\end{bmatrix},
\qquad
B=\begin{bmatrix}
1\\
0
\end{bmatrix},
\qquad
C=\begin{bmatrix}
0 & 1
\end{bmatrix},
\qquad
D=0.
\]
Use \(
\texttt{step(ss(A,B,C,D))}
\) to plot the step response of the system. Find the open-loop poles and choose desired poles to oscillate more slowly than the open-loop system. Write down  the desired characteristic polynomial, and obtain $K$ by comparing it to the characteristic polynomial of the closed-loop system. 
\end{document}
